\documentclass[12pt,a4paper]{article}
\usepackage[T5]{fontenc}
\usepackage[utf8]{inputenc}
\usepackage[vietnamese]{babel}
\usepackage{amsmath,amssymb}
\usepackage{enumitem}
\usepackage{tikz}
\usepackage{pgfplots}
\pgfplotsset{compat=1.18}
\usepackage[margin=2cm]{geometry}
\begin{document}
\begin{center}
Lời giải bài năng lượng liên kết và phân rã hạt nhân
\end{center}

\begin{enumerate}[leftmargin=*,itemsep=1em]

\item
Phản ứng phân rã alpha có dạng
$$
{}^{A}X\longrightarrow {}^{A-4}Y+\alpha.
$$

Năng lượng tỏa ra trong phản ứng là
$$
Q_\alpha
=\left(M_X-M_Y-M_\alpha\right)c^2
=B_Y+B_\alpha-B_X.
$$

Điều kiện để phân rã alpha tự xảy ra là
$$
Q_\alpha>0.
$$

Từ đoạn gần tuyến tính của đồ thị với
$$
A\geq 100,
$$
ta đọc được gần đúng
$$
\frac{B}{A}=a+bA,
\qquad
a=9.6,
\qquad
b=-0.008,
$$
trong đó năng lượng được tính theo $\mathrm{MeV}$.

Do đó
$$
B_X=A(a+bA)
$$
và
$$
B_Y=(A-4)\left[a+b(A-4)\right].
$$

Suy ra
$$
\begin{aligned}
Q_\alpha
&=(A-4)\left[a+b(A-4)\right]+B_\alpha-A(a+bA)\\
&=-4a-8bA+16b+B_\alpha.
\end{aligned}
$$

Với
$$
B_\alpha=25\,\mathrm{MeV},
$$
ta được
$$
Q_\alpha=0.064A-13.528\,\mathrm{MeV}.
$$

Điều kiện
$$
Q_\alpha>0
$$
cho
$$
A>211.375.
$$

Vì số khối là số nguyên nên giá trị ước lượng là
$$
A_a\simeq 212.
$$

\item
\begin{enumerate}[label=\alph*.,leftmargin=2em,itemsep=1em]

\item
Bỏ qua số hạng ghép cặp, với
$$
N=A-Z,
$$
công thức Weizs\"acker trở thành
$$
B=a_vA-a_sA^{2/3}-a_cZ^2A^{-1/3}
-a_a\frac{(A-2Z)^2}{A}.
$$

Với một giá trị $A$ xác định, hạt nhân bền nhất tương ứng với giá trị $Z$ làm
cho $B$ cực đại. Do đó
$$
\frac{\partial B}{\partial Z}
=-2a_cZA^{-1/3}
+\frac{4a_a}{A}(A-2Z)=0.
$$

Suy ra
$$
Z\left(2a_cA^{-1/3}+\frac{8a_a}{A}\right)=4a_a,
$$
hay
$$
Z_{\max}
=\frac{4a_aA}{2a_cA^{2/3}+8a_a}.
$$

Vậy
$$
Z_{\max}
=\frac{A}{2}
\left(1+\frac{a_c}{4a_a}A^{2/3}\right)^{-1}.
$$

\item
Với
$$
A=200,
$$
công thức liên tục ở trên cho
$$
Z_{\max}=79.24.
$$

Vì $A$ chẵn nên $N$ và $Z$ có cùng tính chẵn lẻ. Độ lớn của số hạng ghép cặp là
$$
a_pA^{-3/4}
=33.5\times 200^{-3/4}
=0.630\,\mathrm{MeV}.
$$

So sánh các số nguyên gần $Z_{\max}$ bằng công thức đầy đủ:
$$
\begin{array}{c|c|c}
Z&\delta\ (\mathrm{MeV})&B/A\ (\mathrm{MeV})\\ \hline
78&-0.630&8.0477\\
79&+0.630&8.0458\\
80&-0.630&8.0506
\end{array}
$$

Do đó hạt nhân có năng lượng liên kết riêng lớn nhất ứng với
$$
Z=80.
$$

\item
Với
$$
A=128,
$$
công thức ở câu a cho
$$
Z_{\max}=53.58.
$$

Để so sánh các đồng khối, ta tách phần phụ thuộc vào $Z$ trong năng lượng khối
lượng thành
$$
X(Z)=a_cZ^2A^{-1/3}
+a_a\frac{(A-2Z)^2}{A}+\delta.
$$

Các số hạng còn lại giống nhau đối với mọi hạt nhân có cùng $A$. Tính từ công
thức Weizs\"acker ta được
$$
\begin{array}{c|c}
\text{Hạt nhân}&X\ (\mathrm{MeV})\\ \hline
{}_{51}^{128}\mathrm{Sb}&496.58\\
{}_{52}^{128}\mathrm{Te}&491.18\\
{}_{53}^{128}\mathrm{I}&491.05\\
{}_{54}^{128}\mathrm{Xe}&489.15\\
{}_{55}^{128}\mathrm{Cs}&492.53\\
{}_{56}^{128}\mathrm{Ba}&494.15\\
{}_{57}^{128}\mathrm{La}&501.04
\end{array}
$$

Gọi
$$
\Delta X=X_{\text{mẹ}}-X_{\text{con}}.
$$

Điều kiện năng lượng của các quá trình là
$$
\begin{array}{c|c|c}
\text{Quá trình}&\text{Biến đổi}&\text{Điều kiện}\\ \hline
\beta^-&Z\to Z+1&
\Delta X>m_pc^2+m_ec^2-m_nc^2=-0.79\,\mathrm{MeV}\\
\beta^+&Z\to Z-1&
\Delta X>m_nc^2+m_ec^2-m_pc^2=1.81\,\mathrm{MeV}\\
\text{Bắt electron}&Z\to Z-1&
\Delta X>m_nc^2-m_pc^2-m_ec^2=0.79\,\mathrm{MeV}\\
2\beta^-&Z\to Z+2&
\Delta X>2(m_pc^2+m_ec^2-m_nc^2)=-1.58\,\mathrm{MeV}
\end{array}
$$

Ví dụ, đối với hạt nhân
$$
{}_{53}^{128}\mathrm I,
$$
ta có
$$
\Delta X_{\beta^-}=491.05-489.15=1.90\,\mathrm{MeV}>-0.79\,\mathrm{MeV},
$$
nên phân rã $\beta^-$ xảy ra; đồng thời
$$
\Delta X_{2\beta^-}=491.05-492.53=-1.48\,\mathrm{MeV}>-1.58\,\mathrm{MeV},
$$
nên phân rã hai beta trừ cũng có thể xảy ra về mặt năng lượng.

Kết quả cho ba hạt nhân trong đề là
$$
\begin{array}{c|c|c|c|c}
\text{Hạt nhân}&\beta^-&\beta^+&\text{Bắt electron}&2\beta^-\\ \hline
{}_{53}^{128}\mathrm I&\text{Có}&\text{Không}&\text{Không}&\text{Có}\\
{}_{54}^{128}\mathrm{Xe}&\text{Không}&\text{Không}&\text{Không}&\text{Không}\\
{}_{55}^{128}\mathrm{Cs}&\text{Không}&\text{Có}&\text{Có}&\text{Không}
\end{array}
$$

Trong ba hạt nhân đã cho, hạt nhân bền nhất là
$$
{}_{54}^{128}\mathrm{Xe}.
$$

Cột cuối của bảng được hiểu là phân rã hai beta trừ $2\beta^-$, đúng với mô tả
hai electron được phát ra trong đề.

\end{enumerate}
\end{enumerate}

\newpage
\begin{center}
Lời giải bài mẫu giọt hạt nhân
\end{center}

\begin{enumerate}[leftmargin=*,itemsep=1em]

\item
Mật độ điện tích của quả cầu là
$$
\rho_Q=\frac{Q}{\frac43\pi R^3}
=\frac{3Q}{4\pi R^3}.
$$

Áp dụng định luật Gauss, điện trường bên trong và bên ngoài quả cầu lần lượt là
$$
E(r)=
\begin{cases}
\displaystyle \frac{Qr}{4\pi\varepsilon_0R^3},&r\leq R,\\[2mm]
\displaystyle \frac{Q}{4\pi\varepsilon_0r^2},&r>R.
\end{cases}
$$

Năng lượng tĩnh điện bằng năng lượng của điện trường:
$$
E_c=\frac{\varepsilon_0}{2}\int_{0}^{\infty}E^2\,dV.
$$

Do đó
$$
\begin{aligned}
E_c
&=\frac{\varepsilon_0}{2}
\left[
\int_0^R
\left(\frac{Qr}{4\pi\varepsilon_0R^3}\right)^2
4\pi r^2\,dr
+\int_R^\infty
\left(\frac{Q}{4\pi\varepsilon_0r^2}\right)^2
4\pi r^2\,dr
\right]\\
&=\frac{Q^2}{8\pi\varepsilon_0}
\left(\frac{1}{5R}+\frac{1}{R}\right).
\end{aligned}
$$

Vậy
$$
E_c=\frac{3Q^2}{20\pi\varepsilon_0R}.
$$

\item
Với
$$
Q=Ze,
\qquad
R=R_0A^{1/3},
$$
năng lượng Coulomb là
$$
E_c
=\frac{3Z^2e^2}{20\pi\varepsilon_0R_0A^{1/3}}.
$$

So sánh với số hạng Coulomb trong công thức Weizs\"acker,
$$
E_c=a_3\frac{Z^2}{A^{1/3}},
$$
ta có
$$
R_0=\frac{3e^2}{20\pi\varepsilon_0a_3}.
$$

Đổi đơn vị
$$
a_3=0.72\,\mathrm{MeV}
=0.72\times1.602\times10^{-13}\,\mathrm J,
$$
và dùng
$$
e=1.602\times10^{-19}\,\mathrm C,
$$
ta được
$$
R_0=1.20\times10^{-15}\,\mathrm m.
$$

\item
Khối lượng hạt nhân gần đúng bằng
$$
M\simeq Am.
$$

Vì
$$
R^3=R_0^3A,
$$
mật độ khối lượng của hạt nhân là
$$
\rho_0
=\frac{Am}{\frac43\pi R^3}
=\frac{3m}{4\pi R_0^3}.
$$

Thay số:
$$
\rho_0
=2.3\times10^{17}\,\mathrm{kg\,m^{-3}}.
$$

\item
Năng lượng bề mặt của hạt nhân là
$$
E_s=\sigma S=4\pi\sigma R^2.
$$

Mặt khác, theo công thức Weizs\"acker,
$$
E_s=a_2A^{2/3}.
$$

Với
$$
R^2=R_0^2A^{2/3},
$$
ta được
$$
\sigma=\frac{a_2}{4\pi R_0^2}.
$$

Đổi
$$
a_2=16.8\,\mathrm{MeV}
=16.8\times1.602\times10^{-13}\,\mathrm J,
$$
suy ra
$$
\sigma=1.5\times10^{17}\,\mathrm{N\,m^{-1}}.
$$

\item
Sau phân hạch, hai mảnh có
$$
A_1=kA,
\qquad
Z_1=kZ,
$$
và
$$
A_2=(1-k)A,
\qquad
Z_2=(1-k)Z.
$$

Do tỉ số neutron và proton không đổi, số hạng thể tích và số hạng bất đối xứng
không thay đổi. Độ biến thiên năng lượng bề mặt là
$$
\Delta E_s
=a_2A^{2/3}
\left[k^{2/3}+(1-k)^{2/3}-1\right].
$$

Độ biến thiên năng lượng Coulomb là
$$
\Delta E_c
=a_3\frac{Z^2}{A^{1/3}}
\left[k^{5/3}+(1-k)^{5/3}-1\right].
$$

Phân hạch tự phát thuận lợi về năng lượng khi
$$
\Delta E_s+\Delta E_c<0.
$$

Suy ra điều kiện
$$
\frac{Z^2}{A}>f(k),
$$
trong đó
$$
f(k)
=\frac{a_2}{a_3}
\frac{k^{2/3}+(1-k)^{2/3}-1}
{1-k^{5/3}-(1-k)^{5/3}}.
$$

Hàm số đối xứng qua
$$
k=\frac12
$$
và tăng rất nhanh khi $k$ tiến tới $0$ hoặc $1$.

\begin{center}
\begin{tikzpicture}
\begin{axis}[
    width=11cm,
    height=7cm,
    xmin=0,xmax=1,
    ymin=15,ymax=42,
    xlabel={$k$},
    ylabel={$f(k)$},
    xtick={0,0.2,0.4,0.5,0.6,0.8,1},
    ytick={15,20,25,30,35,40},
    axis line style={black},
    tick style={black},
    grid=major,
    grid style={line width=0.2pt},
    every axis plot/.append style={black},
    clip=true
]
\addplot[thick,domain=0.006:0.994,samples=400]
{(16.8/0.72)*
 (x^(2/3)+(1-x)^(2/3)-1)/
 (1-x^(5/3)-(1-x)^(5/3))};
\addplot[dashed] coordinates {(0.5,15) (0.5,16.3897)};
\addplot[only marks,mark=*] coordinates {(0.5,16.3897)};
\node[above right] at (axis cs:0.5,16.3897)
{$\left(0.5,16.39\right)$};
\end{axis}
\end{tikzpicture}
\end{center}

\item
Do tính đối xứng, giá trị nhỏ nhất đạt tại
$$
k=\frac12.
$$

Khi đó
$$
\begin{aligned}
f_{\min}
&=\frac{a_2}{a_3}
\frac{2(1/2)^{2/3}-1}
{1-2(1/2)^{5/3}}\\
&=16.39.
\end{aligned}
$$

Với hai chữ số có nghĩa,
$$
\left(\frac{Z^2}{A}\right)_{\min}\simeq 16.
$$

\item
Vì elipsoid quay có hai bán trục bằng nhau nên thể tích đúng là
$$
V=\frac43\pi ab^2.
$$

Điều kiện thể tích không đổi cho
$$
ab^2=R^3.
$$

Với
$$
a=R(1+\varepsilon),
$$
ta được
$$
b=\frac{R}{\sqrt{1+\varepsilon}}
\simeq R\left(1-\frac{\varepsilon}{2}\right).
$$

Nếu đặt
$$
b=R\lambda,
$$
thì
$$
\lambda=\frac{1}{\sqrt{1+\varepsilon}}
\simeq1-\frac{\varepsilon}{2}.
$$

Nếu giữ đúng ký hiệu $b=R/\lambda$ trong bản đề thì hệ số tương ứng là nghịch
đảo:
$$
\lambda=\sqrt{1+\varepsilon}
\simeq1+\frac{\varepsilon}{2}.
$$

Quan hệ hình học không phụ thuộc cách ký hiệu vẫn là
$$
b=\frac{R}{\sqrt{1+\varepsilon}}.
$$

\begin{center}
\begin{tikzpicture}[>=stealth,scale=1.0]
  \draw (0,0) ellipse (3.2cm and 1.75cm);
  \draw[->] (0,0)--(3.2,0) node[midway,below] {$a$};
  \draw[->] (0,0)--(0,1.75) node[midway,left] {$b$};
  \draw[dashed] (-3.6,0)--(3.6,0);
  \draw[dashed] (0,-2.1)--(0,2.1);
  \fill (0,0) circle (1.5pt);
\end{tikzpicture}
\end{center}

\item
Với
$$
a=R(1+\varepsilon),
\qquad
b=c=\frac{R}{\sqrt{1+\varepsilon}},
$$
công thức gần đúng của diện tích mặt elipsoid cho
$$
S
=4\pi R^2
\left[
\frac{2(1+\varepsilon)^{0.8}
+(1+\varepsilon)^{-1.6}}{3}
\right]^{0.625}.
$$

Khai triển đến bậc hai:
$$
(1+\varepsilon)^{0.8}
\simeq1+0.8\varepsilon-0.08\varepsilon^2
$$
và
$$
(1+\varepsilon)^{-1.6}
\simeq1-1.6\varepsilon+2.08\varepsilon^2.
$$

Do đó
$$
\begin{aligned}
S
&\simeq4\pi R^2
\left(1+0.64\varepsilon^2\right)^{0.625}\\
&\simeq4\pi R^2
\left(1+0.4\varepsilon^2\right).
\end{aligned}
$$

Độ tăng diện tích mặt là
$$
\Delta S
=\frac25\varepsilon^2\,4\pi R^2.
$$

Vì năng lượng bề mặt ban đầu bằng
$$
4\pi\sigma R^2=a_2A^{2/3},
$$
nên độ tăng năng lượng bề mặt là
$$
\Delta E_s
=\frac25a_2A^{2/3}\varepsilon^2.
$$

Theo dữ kiện của đề, độ biến thiên năng lượng Coulomb là
$$
\Delta E_c
=-\frac15E_c\varepsilon^2
=-\frac15a_3\frac{Z^2}{A^{1/3}}\varepsilon^2.
$$

Trạng thái cầu mất ổn định khi tổng năng lượng không tăng dưới một biến dạng
nhỏ:
$$
\Delta E_s+\Delta E_c\leq0.
$$

Tại giá trị tới hạn,
$$
\frac25a_2A^{2/3}
=\frac15a_3\frac{Z^2}{A^{1/3}}.
$$

Vì vậy
$$
\left(\frac{Z^2}{A}\right)_{\text{tới hạn}}
=\frac{2a_2}{a_3}
=\frac{2\times16.8}{0.72}
=46.7.
$$

\end{enumerate}
\end{document}
